\documentclass[../proyecto.tex]{subfiles}
\graphicspath{{\subfix{../images/}}}
\begin{document}

\section{Formatos}

% máximo 1 página

% Indicar cómo se espera abordar la tesis a nivel de escritura, y otras dimensiones extratextuales

Esta tesis incluye a grandes rasgos dos secciones fundamentales. La primera es el documento textual de la tesis doctoral que recopila y establece los fundamentos teóricos y conceptuales de la investigación realizada en electrónica y computación aplicada a popusintesíntesis. La segunda es su materialización con una colección de publicaciones disciplinares propias de la electrónica y la computación, que permiten el entendimiento de los procesos de diseño, fabricación y documentación de los artefactos técnicos de la tesis.

En la tesis se incluyen capítulos donde se proponen y describen módulos necesarios y opcionales, suficientes e insuficientes para crear sintetizadores sonoras, incluyendo metáforas conceptuales, contexto histórico, contexto cultural, estrategias de producción artesanal, estrategias de organización empresarial para compras masivas y ventas menos masivas de materiales para artes, y estrategias académicas de pasantes y practicantes que han expandido el trabajo en este proyecto a lo largo de los años.

A nivel de acompañamiento electrónico, se incluyen diagramas de circuitos, listas de materiales, estrategias de diseño y fabricación de placas impresas, y los manuales técnicos de uso y reparación de sintetizadores electrónicos.

A nivel de acompañamiento computacional, se incluyen diagramas de flujo, pseudocódigo, y código fuente desarrollado con control de versiones. Este código es usado tanto para la síntesis de sonido dentro los sintetizadores realizados, como para la construcción de carcasas y otros artefactos físicos necesarios para esta misión.

\end{document}
